\chapter{NUMERICAL RESULTS}
\label{chap:results}

This chapter reports and interprets the experiments described in Chapter~\ref{chap:methodology}. The presentation is organized around a single question: how faithfully does each architecture, in each training regime, reproduce the pseudospectral reference, and how does that fidelity depend on the wavenumber, on the elapsed time, on the nonlinearity, and on the discretization? Section~\ref{sec:res_training} reports the training cost. Section~\ref{sec:res_ntk} examines the Neural Tangent Kernel diagnostic that grounds the discussion of spectral bias. Section~\ref{sec:res_spectral_bias} follows the frequency content of the error along training, and Section~\ref{sec:res_spectral_fidelity} inspects the final spectral fidelity at the intermediate parameter. Section~\ref{sec:res_gibbs} isolates the temporal-boundary artifact that separates the data-driven operator from its physics-informed variant, Section~\ref{sec:res_nonlinearity} sweeps the nonlinearity axis, and Section~\ref{sec:res_resolution} tests the zero-shot resolution behaviour. Section~\ref{sec:res_synthesis} draws the threads together. Throughout, the pseudospectral solver is treated as ground truth, and no architecture is credited with a success it does not demonstrate: every method examined here leaves a clearly identifiable failure mode.

% ===============================================================
\section{Training Cost}
\label{sec:res_training}

Table~\ref{tab:training} summarizes the size, wall-clock training time and final training loss of the five trained models. The two operators are roughly five times larger than the single-parameter PINN and take about a fifth longer to train, but in exchange a single operator covers the entire $20$-point parameter family~\eqref{eq:param_grid}, whereas each PINN is bound to the single parameter $\alpha = \beta = 3.21$ and would have to be retrained from scratch for any other value. Amortized over the family, the operators are by far the cheaper way to obtain a solution at a new parameter, since inference is a single forward pass.

\begin{table}[htbp]
    \centering
    \caption{Model size, training time on the GTX~1660~SUPER, and final training loss. The final-loss column is \emph{not} comparable across rows: each objective is a different combination of data, residual and initial-condition terms, so its magnitude reflects the composition of the loss rather than solution accuracy. Solution accuracy is measured instead by the relative-error metrics of the following sections.}
    \label{tab:training}
    \small
    \begin{tabular}{llrrr}
        \hline
        \textbf{Model} & \textbf{Regime} & \textbf{Parameters} & \textbf{Train time (s)} & \textbf{Final loss} \\
        \hline
        FNO  & supervised   & $1{,}057{,}698$ & $1121.6$ & $5.9 \times 10^{-6}$ \\
        PINO & with data    & $1{,}057{,}698$ & $1175.5$ & $2.2 \times 10^{-5}$ \\
        PINO & no data      & $1{,}057{,}698$ & $1176.2$ & $7.2 \times 10^{-6}$ \\
        PINN & with data    & $198{,}658$     & $861.8$  & $1.1 \times 10^{-4}$ \\
        PINN & no data      & $198{,}658$     & $842.2$  & $6.4 \times 10^{-6}$ \\
        \hline
    \end{tabular}
    \\[0.4em]
    \autoriapropria
\end{table}

A caution is worth stating at the outset. The final-loss column tempts a ranking, but the losses are not on a common scale: the PINN with data reaches a larger final loss ($1.1 \times 10^{-4}$) than the PINN without data ($6.4 \times 10^{-6}$) simply because the added supervised term contributes its own residual to the total, not because it is a worse solution --- indeed the following sections show the opposite. Training loss is a poor proxy for accuracy here, and all comparisons below rest on the held-out relative-error metrics of Section~\ref{subsec:met_defs}.

% ===============================================================
\section{The Mechanism of Spectral Bias: NTK Diagnostic}
\label{sec:res_ntk}

Before comparing architectures, we verify that the spectral-bias mechanism derived in Section~\ref{subsec:spectral_bias} is actually operative for the Boussinesq PINN. Figure~\ref{fig:ntk_panel} reports the Neural Tangent Kernel diagnostic of Section~\ref{sec:met_ntk} for three network widths.

\begin{figure}[htbp]
    \centering
    \begin{subfigure}[t]{0.48\textwidth}
        \centering
        \includegraphics[width=\linewidth]{pinn/ntk/pinn_ntk_panel_param.png}
        \caption{Relative parameter drift $\|\theta-\theta_0\|/\|\theta_0\|$.}
        \label{fig:ntk_param}
    \end{subfigure}
    \hfill
    \begin{subfigure}[t]{0.48\textwidth}
        \centering
        \includegraphics[width=\linewidth]{pinn/ntk/pinn_ntk_panel_kernel.png}
        \caption{Relative kernel drift $\|K-K_0\|/\|K_0\|$.}
        \label{fig:ntk_kernel}
    \end{subfigure}

    \vspace{0.8em}
    \begin{subfigure}[t]{0.48\textwidth}
        \centering
        \includegraphics[width=\linewidth]{pinn/ntk/pinn_ntk_panel_spectrum.png}
        \caption{Eigenvalue spectrum of $K$: initialization (solid) versus end of training (dashed), on log--log axes.}
        \label{fig:ntk_spectrum}
    \end{subfigure}
    \caption{Neural Tangent Kernel diagnostic for the Boussinesq PINN residual at $\alpha = \beta = 3.21$, for widths $8$, $128$ and $512$.}
    \label{fig:ntk_panel}
    \autoriapropria
\end{figure}

Two predictions of the theory are confirmed. First, the frozen-kernel assumption underlying Equation~\eqref{eq:error_linear_ode} becomes more accurate as the width grows: the widest network ($512$) drifts by less than $0.1\%$ in its parameters and by under $10\%$ in its kernel over the whole run, whereas the narrowest ($8$) drifts by several percent in parameters and by more than $40\%$ in its kernel. Wide networks sit closer to the lazy-training regime in which the analysis is exact, and the ordering of the curves is monotone in width. Second, and more importantly, the eigenvalue spectrum in the right panel spans many orders of magnitude and decays precipitously: only a handful of eigenvalues carry appreciable weight before the spectrum collapses toward the numerical floor. By Equation~\eqref{eq:training_time}, the time to resolve an eigendirection is inversely proportional to its eigenvalue, so this steep decay means that all but the leading, smooth eigendirections are learned prohibitively slowly. This is precisely the eigenvalue disparity that \textcite{wang2022ntk} identify as the source of spectral bias, and it sets the expectation for everything that follows: the PINN should acquire the low-frequency structure of the solution readily and stall on the high-frequency detail.

% ===============================================================
\section{Spectral Bias Along Training}
\label{sec:res_spectral_bias}

Figure~\ref{fig:spectral_bias} tracks the relative spectral error in the low, mid and high wavenumber bands of Equation~\eqref{eq:bands} as training proceeds, for all five models. Read within each panel, the figure tells a consistent story; read across panels, it must be interpreted with the caveat of Section~\ref{subsec:met_defs}, namely that the operators are monitored on the near-linear sample $\alpha = \beta = 0.1$ while the PINNs are monitored on the harder $\alpha = \beta = 3.21$. Cross-family magnitudes are therefore not directly comparable in this figure; the apples-to-apples cross-family comparison is deferred to Section~\ref{sec:res_spectral_fidelity}, where every model is evaluated at the common parameter $3.21$.

\begin{figure}[htbp]
    \centering
    \includegraphics[width=\textwidth]{spectral_bias_evolution.png}
    \caption{Relative spectral error by frequency band during training. In every panel the low band (blue) is learned first and to the lowest error, the mid band (purple) trails it, and the high band (red) is learned last and least --- the signature of spectral bias. Note the mixed reference samples: operators are tracked at $\alpha = \beta = 0.1$, PINNs at $\alpha = \beta = 3.21$.}
    \label{fig:spectral_bias}
    \autoriapropria
\end{figure}

The within-panel ordering is unambiguous and universal: in all five cases the low band is driven lowest, the mid band trails, and the high band is learned last and least. \emph{No architecture and no regime escapes this ordering}, which is the empirical face of the NTK spectrum of Figure~\ref{fig:ntk_panel}. What differs is the magnitude of the gap and how far each band is driven down.

The comparisons that \emph{are} valid in this figure --- same architecture, same reference sample, differing only in the data term --- are decisive. For the PINN at $3.21$, adding data lowers the terminal low-band error from about $0.16$ to about $0.06$ and the mid-band error from about $0.33$ to about $0.18$, while leaving the high band essentially untouched near $0.70$. For the PINO at $0.1$, the effect is far larger: the data term drives the low band from about $0.27$ down to below $10^{-2}$ and the mid band from about $0.62$ to about $0.02$, again with a much smaller improvement in the high band, which falls only from about $0.81$ to about $0.45$. In both families the lesson is the same: \textbf{the supervised term is a powerful lever on the low and mid bands and a weak one on the high band.} The physics residual alone, in the two no-data panels, never brings any band close to the levels reached with data; the PINO~(no~data) high band barely moves over the entire run.

The FNO panel, monitored at the same easy parameter as the PINO, reaches the lowest low- and mid-band errors of the whole figure (about $0.013$ and $0.026$) and the lowest high-band error as well (about $0.28$). Even so, its high band sits two orders of magnitude above its low band. The single most honest summary of Figure~\ref{fig:spectral_bias} is therefore not that operators ``solve'' spectral bias but that they \emph{compress} it: the best model still leaves a high-frequency residual of roughly $28\%$, and every other model leaves more.

% ===============================================================
\section{Spectral Fidelity at the Intermediate Regime}
\label{sec:res_spectral_fidelity}

To compare the families fairly, we now fix the parameter at the intermediate value $\alpha = \beta = 3.21$ and inspect the full spatial spectrum of the prediction at three times, $t = 0$, $t = T/2 = 15$ and $t = T = 30$. Figures~\ref{fig:spec_pinn_nodata}--\ref{fig:spec_pino_nodata} report the five models; in every panel the black curve is the reference spatial spectrum and the dashed orange curve is the prediction, with the accompanying relative-error curves and the time-averaged error at the foot of each figure.

\begin{figure}[htbp]
    \centering
    \includegraphics[width=\textwidth]{pinn/no_data/pinn_no_data_model2_spectral_panel.png}
    \caption{PINN (no data) spatial spectrum at $\alpha = \beta = 3.21$. Already at $t = 0$ the predicted spectrum undershoots the reference tail, and the relative error rises monotonically with wavenumber; the mean spectral error grows from $0.17$ at $t=0$ to $0.70$ at $t=30$.}
    \label{fig:spec_pinn_nodata}
    \autoriapropria
\end{figure}

\begin{figure}[htbp]
    \centering
    \includegraphics[width=\textwidth]{pinn/with_data/pinn_model2_spectral_panel.png}
    \caption{PINN (with data) spatial spectrum at $\alpha = \beta = 3.21$. The low-wavenumber crests are tracked more closely than in Figure~\ref{fig:spec_pinn_nodata}, but the high-wavenumber undershoot and the growth of the mean error in time ($0.17 \to 0.58 \to 0.68$) persist.}
    \label{fig:spec_pinn_data}
    \autoriapropria
\end{figure}

The two PINN panels expose the pathology in its clearest form. In the no-data case (Figure~\ref{fig:spec_pinn_nodata}) the deficiency is visible already at $t = 0$: the network cannot even represent the initial $\sech^2$ profile at high wavenumber, and the relative error climbs monotonically toward unity as the mode index increases. By $t = 30$ the predicted spectrum has lost most of its mid- and high-frequency amplitude, and the mean spectral error reaches $0.70$. Adding data (Figure~\ref{fig:spec_pinn_data}) tightens the low-wavenumber crests noticeably --- consistent with the low-band improvement of Figure~\ref{fig:spectral_bias} --- yet the aggregate mean error barely changes ($0.17 \to 0.58 \to 0.68$), because it is dominated by the mid and high modes that the data term fails to recover at this hard parameter. Data helps the PINN where the NTK is favourable and does not rescue it where the NTK is adverse.

\begin{figure}[htbp]
    \centering
    \includegraphics[width=\textwidth]{fno/fno_model2_spectral_panel.png}
    \caption{FNO spatial spectrum at $\alpha = \beta = 3.21$. Low and mid wavenumbers are tracked almost exactly (mean error $0.013$ at $t=0$), but the error migrates to the high modes and grows steadily in time, reaching a mean of $0.40$ at $t=30$.}
    \label{fig:spec_fno}
    \autoriapropria
\end{figure}

\begin{figure}[htbp]
    \centering
    \includegraphics[width=\textwidth]{pino/with_data/pino_model2_spectral_panel.png}
    \caption{PINO (with data) spatial spectrum at $\alpha = \beta = 3.21$. The prediction tracks the reference across the band and, unlike the FNO, the mean error remains bounded in time ($0.077 \to 0.13 \to 0.20$): the physics residual stabilizes the temporal evolution.}
    \label{fig:spec_pino_data}
    \autoriapropria
\end{figure}

\begin{figure}[htbp]
    \centering
    \includegraphics[width=\textwidth]{pino/no_data/pino_no_data_model2_spectral_panel.png}
    \caption{PINO (no data) spatial spectrum at $\alpha = \beta = 3.21$. Without the supervised term the low- and mid-band undershoot returns (mean error $0.11 \to 0.39 \to 0.33$), confirming that the physics residual alone does not overcome spectral bias.}
    \label{fig:spec_pino_nodata}
    \autoriapropria
\end{figure}

The operators, learning directly in the frequency domain, reverse the picture at low and mid wavenumbers. The FNO (Figure~\ref{fig:spec_fno}) reproduces the low and mid spectrum almost exactly at early times --- a mean error of $0.013$ at $t = 0$, an order of magnitude below the PINN --- because those are precisely the modes its spectral layers parametrize. Its weakness is elsewhere: the error concentrates in the high modes and, crucially, \emph{grows monotonically in time}, from a mean of $0.013$ at $t = 0$ to $0.40$ at $t = 30$. The PINO with data (Figure~\ref{fig:spec_pino_data}) inherits the operator's spectral accuracy but does not share this temporal degradation: its mean error rises only from $0.077$ to $0.20$ over the same horizon, remaining the most stable of all five models at late times. The contrast is sharp at $t = 30$, where the PINO with data ($0.20$) is markedly better than the FNO ($0.40$), even though the FNO was better at $t = 0$; the physics residual trades a slightly worse start for a much steadier trajectory. Finally, the PINO without data (Figure~\ref{fig:spec_pino_nodata}) forfeits this advantage: stripped of the supervised term, its low- and mid-band undershoot returns and its mean error climbs to $0.33$--$0.39$, midway between the operators-with-data and the PINNs. The physics constraint, on its own, is not a substitute for data.

% ===============================================================
\section{The Temporal-Boundary Artifact}
\label{sec:res_gibbs}

The temporal growth of the FNO error observed in Figure~\ref{fig:spec_fno} has a specific and instructive cause. As explained in Section~\ref{sec:met_fno}, the operator applies the Fourier transform along the temporal axis as well as the spatial one, implicitly assuming periodicity in time. Because the Boussinesq trajectory is not periodic --- the field at $t = T$ does not match the field at $t = 0$ --- this assumption creates an artificial discontinuity at the temporal boundary and excites the Gibbs phenomenon there \parencite{gibbs1899}. Figures~\ref{fig:ab_fno} and~\ref{fig:ab_pino} make the effect and its remedy visible.

\begin{figure}[htbp]
    \centering
    \includegraphics[width=\textwidth]{fno/fno_model2_alpha_beta_panel.png}
    \caption{FNO across the nonlinearity axis ($\alpha=\beta \in \{0.1, 3.21, 4.2\}$), evaluated at resolution $256$. The predicted surfaces (left) capture the two counter-propagating pulses, but the time-resolved relative error (right) exhibits a pronounced spike at the final time $t = 30$: the temporal Gibbs artifact from treating time as a periodic axis.}
    \label{fig:ab_fno}
    \autoriapropria
\end{figure}

\begin{figure}[htbp]
    \centering
    \includegraphics[width=\textwidth]{pino/with_data/pino_model2_alpha_beta_panel.png}
    \caption{PINO (with data) across the same nonlinearity axis and resolution. The predicted surfaces are visually indistinguishable from the FNO, but the terminal-time spike of Figure~\ref{fig:ab_fno} is absent: the error rises only gently toward $t = 30$. The physics residual and initial-condition term act as a temporal regularizer.}
    \label{fig:ab_pino}
    \autoriapropria
\end{figure}

In the FNO panel (Figure~\ref{fig:ab_fno}) every parameter value shows the same feature: after an initial transient the time-resolved error settles to a plateau and then jumps sharply at $t = 30$, reaching values two to three times the plateau level. This end-time spike is present at all three nonlinearities and is the dominant contribution to the FNO's late-time error. In the PINO panel (Figure~\ref{fig:ab_pino}) the predicted surfaces are visually identical, yet the spike is gone: the error curve rises only slightly toward the boundary. The explanation is that the PDE residual and the initial-condition term penalize exactly the temporal inconsistency that the periodic transform would otherwise tolerate, damping the boundary oscillation. This confirms, and explains, an observation that motivated the experiment: the Gibbs artifact of the FNO is suppressed once the physics is added. It is, however, a suppression and not an elimination --- a small terminal rise remains --- and it comes at the modelling cost of the residual machinery.

% ===============================================================
\section{Accuracy Across the Nonlinearity Axis}
\label{sec:res_nonlinearity}

The box plots at the foot of Figures~\ref{fig:ab_fno} and~\ref{fig:ab_pino} summarize the distribution of the time-resolved error~\eqref{eq:rel_time} over the whole trajectory, at each of the three parameters. Their most striking feature is counterintuitive: the \emph{smallest} nonlinearity, $\alpha = \beta = 0.1$, carries the \emph{largest} relative error for both operators --- a median near $0.18$, against roughly $0.12$--$0.13$ at the more nonlinear values $3.21$ and $4.2$. The strongly nonlinear regime, which one might expect to be hardest, is in fact reproduced with the lowest relative error.

We resist a triumphant reading of this and offer a measured one. Two effects plausibly combine. First, $\alpha = \beta = 0.1$ sits at the extreme lower edge of the training grid~\eqref{eq:param_grid}, where the operator has the least surrounding data to interpolate from, so a boundary-of-range penalty is expected. Second, and more mechanically, the relative metric~\eqref{eq:rel_time} normalizes by the instantaneous norm of the reference: in the near-linear regime the solution is dominated by the two clean, slowly deforming pulses of the wave-equation limit, whose energy is concentrated and whose norm is modest, so a fixed absolute error translates into a larger relative one. The nonlinear cases spread more energy across the domain and are, in relative terms, more forgiving. What the box plots do \emph{not} support is the claim that these operators become uniformly more accurate as nonlinearity grows in any absolute sense; they show only that the relative error, under this normalization and this training grid, is non-monotone in the parameter, and that the near-linear boundary case is the least well served.

% ===============================================================
\section{Zero-Shot Resolution Behaviour}
\label{sec:res_resolution}

A central promise of neural operators is discretization-invariance: an operator trained at one resolution can be queried at another without retraining. Figure~\ref{fig:res_pino} tests this for the PINO with data, trained at resolution $128$ and evaluated at $128$, $256$ and $512$ at the fixed parameter $\alpha = \beta = 3.21$.

\begin{figure}[htbp]
    \centering
    \includegraphics[width=\textwidth]{pino/with_data/pino_model2_resolution_panel.png}
    \caption{PINO (with data) queried zero-shot at resolutions $128$, $256$ and $512$ at $\alpha = \beta = 3.21$. The operator produces a coherent field at every resolution, but the relative error against the pseudospectral reference \emph{grows} with resolution (median error rising from about $0.07$ at $128$ to about $0.16$ at $512$): discretization-invariance of form is not accuracy-invariance.}
    \label{fig:res_pino}
    \autoriapropria
\end{figure}

The qualitative promise holds: at every resolution the operator returns a smooth, coherent surface with the correct pulse structure, and it does so without any retraining or architectural change --- the input is simply re-sampled on the finer grid. The quantitative promise is more nuanced. The box plot shows the relative error \emph{increasing} with resolution, from a median near $0.07$ at the training resolution of $128$ to about $0.12$ at $256$ and $0.16$ at $512$. The reason is that the finer references contain genuine high-wavenumber content that the operator, trained at $128$ and retaining only $16$ Fourier modes per axis, never learned to produce; querying it against a sharper reference simply exposes the fine scales it is missing. There is also a secondary effect at the training resolution itself, where the error oscillates in time in step with the wave reflections, an aliasing signature of the coarse grid. The honest conclusion is that zero-shot superresolution is real as a \emph{geometric} property --- the operator accepts and returns fields at any resolution --- but not as an \emph{accuracy} property: fidelity is anchored to the spectral band seen in training, and evaluating on a richer grid can only reveal, not recover, the modes left out.

% ===============================================================
\section{Synthesis}
\label{sec:res_synthesis}

Taken together, the experiments resist any single-winner narrative and instead map out where each method succeeds and where it fails.

\begin{itemize}
    \item \textbf{Spectral bias is universal.} Every architecture, in every regime, learns low wavenumbers first and best and high wavenumbers last and least (Figures~\ref{fig:ntk_panel}, \ref{fig:spectral_bias}). No method drives the high band below roughly $0.28$ relative error; the best in that band is the plain FNO.

    \item \textbf{Data is the decisive lever, physics is not a substitute for it.} Within both the PINN and the PINO families, adding the supervised term sharply reduces the low- and mid-band error, while the physics-only regimes revert toward the biased baseline (Figures~\ref{fig:spectral_bias}, \ref{fig:spec_pino_data}, \ref{fig:spec_pino_nodata}). Physics alone neither matches data nor removes the bias.

    \item \textbf{Operators dominate at low and mid wavenumbers; PINNs do not.} At the common parameter $3.21$, the FNO and the PINO with data track the low and mid spectrum an order of magnitude more accurately than either PINN (Section~\ref{sec:res_spectral_fidelity}), and cover the whole parameter family from one training, whereas each PINN is bound to a single parameter.

    \item \textbf{Physics buys temporal stability.} The FNO's error grows in time and spikes at the temporal boundary through the Gibbs artifact of its periodic time axis; the PINO's residual suppresses that spike and keeps the late-time error bounded (Section~\ref{sec:res_gibbs}). This is the clearest concrete benefit of the physics term in this study.

    \item \textbf{The remaining limitations are real.} The near-linear boundary parameter is the least accurate (Section~\ref{sec:res_nonlinearity}); zero-shot superresolution degrades in accuracy as the query grid sharpens (Section~\ref{sec:res_resolution}); and the high-frequency band is never resolved by any method. These are not artefacts of a single run but consistent features of the design.
\end{itemize}

The picture that emerges is the one anticipated in Chapter~\ref{chap:intro}: these methods complement, rather than replace, the pseudospectral solver that defines the accuracy ceiling here. Their value lies in amortizing the cost of the parameter sweep and in the temporal regularization that physics supplies, not in a uniform accuracy advantage --- and least of all at the high wavenumbers where the reference solver remains unchallenged.
